\section{Discussão dos resultados}
\label{sec:cap4_app_final_discussao}

Os resultados apresentados neste capítulo avaliam a cadeia de aproximações da missão de \gls{rvdb}, desde a transferência orbital inicial até a aproximação final.
Em resumo, buscou-se verificar se:
(i) a transferência de Hohmann conduz o perseguidor a uma órbita adequada para a transição ao referencial \gls{lvlh};
(ii) o controle translacional baseado nas equações \gls{hcw} leva o perseguidor à região de aproximação de curta distância com erro e velocidade controlados;
(iii) o controle de atitude estabiliza o perseguidor tanto com o manipulador inoperante quanto em regime acoplado base--manipulador;
e (iv) o manipulador posiciona a ferramenta em uma vizinhança próxima da porta de acoplamento, sem comprometer a dinâmica global do sistema.

\subsection{Coerência entre as fases de aproximação}

A linha do tempo da missão mostra que as três aproximações (longa, curta e final) são encadeadas de maneira consistente.
Na aproximação de longa distância, a transferência de Hohmann eleva o raio orbital do perseguidor da altitude inicial de $500$ km para aproximadamente $599{,}9$ km, praticamente coincidente com a órbita circular do alvo, que se encontra a $600$ km de altitude. A diferença radial residual após a segunda queima é da ordem de $100$ m, valor coerente com o emprego posterior de um modelo linearizado em LVLH para descrever o movimento relativo.

A análise da longitude inercial mostra que a condição de fase entre perseguidor e alvo é ajustada;
ao final da transferência e do \emph{drift} orbital, a diferença de longitude reduz-se para próxmo de $0{,}0041^\circ$, o que corresponde a uma separação relativa de aproximadamente $506$ m. 

A comparação entre os cenários (A) e (B) indica que alterações nas condições iniciais (como altitude e longitude do perseguidor) preservam a estrutura relativa da missão, mas impactam de forma sensível o tempo total e o consumo de $\Delta v$.
No cenário B, por exemplo, o tempo até o término da aproximação de longa distância aumenta em cerca de $365\%$, e o consumo total de $\Delta v$ praticamente dobra em relação ao cenário de referência. Isso confirma que a escolha da órbita inicial e da fase inercial desempenha papel relevante no custo propulsivo da missão, ainda que o estado final para o início do movimento relativo em LVLH permaneça semelhante.

\subsection{Desempenho do controle translacional em LVLH}

Na aproximação de curta distância, a dinâmica relativa no referencial LVLH mostra que a lei de controle proporcional--derivativa cumpre os objetivos de redução da posição e da velocidade relativas, mantendo o perseguidor dentro de limites operacionais pré-estabelecidos. As componentes de posição partem de valores da ordem de centenas de metros e convergem para uma condição final próxima de $-0{,}007$ m em along-track e $-0{,}500$ m em radial, enquanto a componente cross-track permanece praticamente nula desde o início da fase.

Os tempos de assentamento em posição, calculados com critério de tolerância de $\pm 0{,}05$ m durante um intervalo mínimo de $15$ s, são da ordem de $515$ s para a componente radial e $530$ s para a componente along-track, com assentamento imediato em cross-track. Esses tempos refletem um compromisso entre rapidez de convergência e limitação de esforços de controle, uma vez que a mesma lei deve respeitar restrições de velocidade relativa e de aceleração impostas para garantir segurança na vizinhança do alvo.

A velocidade relativa inicial é dominada pela componente along-track, cerca de $0{,}1625$ m/s, e converge para valores finais da ordem de $10^{-5}$ m/s. A saturação da velocidade em $0{,}3000$ m/s para distâncias superiores a $20$ m e em $0{,}0300$ m/s na vizinhança do alvo impede que a dinâmica aproximada pelas equações de \gls{hcw} conduza a velocidades de alguns metros por segundo. Os picos de velocidade não saturada (cerca de $4{,}51$ m/s em radial e $8{,}62$ m/s em along-track) mostram que, sem essa limitação, a manobra apresentaria risco de colisão entre o perseguidor e o alvo. Em contrapartida, a aceleração de controle é mantida dentro de $\pm 1{,}0000$ m/s², e tende rapidamente a zero após o regime transitório, o que caracteriza um esforço de controle compatível com atuadores de baixa potência.
Esses resultados indicam que o controle translacional consegue trazer o perseguidor para uma região de referência próxima ao alvo, com limites de posição e velocidade adequados para a transição à fase de aproximação final, ainda que à custa de tempos de assentamento relativamente longos.

\subsection{Controle de atitude e dinâmica da base com manipulador acoplado}

Durante toda a missão, o controle de atitude é responsável por alinhar o perseguidor ao referencial do alvo e por rejeitar as perturbações induzidas pelo manipulador robótico. Os erros iniciais de atitude em relação ao alvo — cerca de $7^\circ$ em roll, $8^\circ$ em pitch e $9^\circ$ em yaw — são reduzidos para valores finais próximos de $0{,}5^\circ$ em cada componente, dentro do limite de $\pm 1^\circ$ estabelecido. O alinhamento se mantém tanto na fase em que a dinâmica é descrita pelas equações de Euler quanto após a comutação para o modelo de Lagrange com base flutuante.

As velocidades angulares da base permanecem em faixas moderadas, com limites aproximados de $[-1{,}20,\;1{,}20]$ rad/s em $w_x$, $[-1{,}26,\;1{,}25]$ rad/s em $w_y$ e $[-0{,}96,\;1{,}00]$ rad/s em $w_z$, convergindo ao final para valores residuais de $0{,}01$ rad/s, $-0{,}01$ rad/s e $0{,}03$ rad/s, respectivamente. As acelerações angulares apresentam picos concentrados na transição entre os modelos dinâmicos, chegando a cerca de $-109{,}9$ rad/s² em $w\dot{y}$, mas retornam a aproximadamente $\pm 2{,}00$ rad/s² no regime final. Esses resultados sugerem que a comutação entre os modelos produz transitórios mais intensos, porém de duração limitada, após os quais a base evolui para um regime de rotação lenta, compatível com a fase de engate.

Os torques de controle de atitude (PD) oscilam tipicamente na faixa de poucos N·m, com picos absolutos próximos de $12$ N·m em módulo e valores finais de aproximadamente $-0{,}016$ N·m, $0{,}028$ N·m e $-0{,}58$ N·m em $x$, $y$ e $z$, respectivamente.
A versão compensada desses torques, que inclui termos de acoplamento dinâmico base--manipulador, apresenta valores finais muito próximos, com diferenças da ordem de $10^{-3}$ a $10^{-2}$ N·m, que afetam principalmente a forma transitória. Os torques de reação na base, induzidos pela movimentação do manipulador, permanecem significativamente menores; variam, em módulo, em torno de $0{,}03$ N·m, com valores finais inferiores a $0{,}01$ N·m em cada componente. Assim, o controlador de atitude consegue rejeitar as perturbações associadas ao manipulador sem exigir níveis de torque incompatíveis com uma plataforma de pequeno porte.

A evolução da energia cinética e do momento angular reforça essa interpretação. A energia cinética total atinge valor máximo de cerca de $0{,}1914$ J e decai para aproximadamente $1{,}35\times 10^{-4}$ J ao final da aproximação, enquanto as componentes do momento angular total se mantêm em faixas da ordem de $10^{-1}$ N·m·s, com valores residuais entre $10^{-3}$ e $10^{-2}$ N·m·s. A manobra termina, portanto, em um estado de baixa energia cinética e baixo momento angular residual.

\subsection{Atuação do manipulador e aproximação geométrica da ferramenta}

O manipulador robótico opera durante a aproximação final para posicionar a ferramenta em uma configuração adequada ao engate. As posições articulares convergem para uma configuração bem definida em regime final, aproximadamente $q_1 = -1{,}5820$ rad, $q_2 = -0{,}2649$ rad e $q_3 = 2{,}4860$ rad, com velocidades residuais reduzidas (por exemplo, $\dot{q}_1 \approx 0{,}0097$ rad/s e $\dot{q}_2 \approx 0{,}0009$ rad/s). Os torques de junta permanecem na faixa de décimos de N·m, com valores finais próximos de zero nas juntas $1$ e $2$ e esforço mais significativo na junta $3$, que responde pela extensão principal da ferramenta.

Os erros de seguimento em relação à referência de cinemática inversa permanecem contidos: as faixas de erro de posição por junta são da ordem de $10^{-3}$ a $10^{-2}$ rad, com valores finais próximos de zero em $q_2$ e da ordem de milésimos de radiano em $q_1$ e $q_3$. Os erros de velocidade exibem comportamento similar, o que indica que a trajetória descrita pelo manipulador acompanha de perto a referência planejada, sem oscilações significativas.

Do ponto de vista geométrico, em relação à referência desejada para a posição da ponta, os erros de posição $(e_x,e_y,e_z)$ variam ao longo da aproximação dentro de faixas inferiores a $0{,}12$ m em módulo, com valores finais de cerca de $-0{,}0015$ m, $-0{,}0045$ m e $1{,}61\times 10^{-4}$ m, respectivamente. A norma do erro de posição permanece abaixo de aproximadamente $0{,}93$ m e converge para um valor final de $4{,}7\times 10^{-3}$ m, isto é, poucos milímetros em relação à posição de referência.
Sendo assim, o sistema consegue posicionar o efetuador em uma região compatível com uma janela de captura fina. Entretanto, o modelo não impõe explicitamente uma restrição de alinhamento angular perfeito entre a ferramenta e a porta; a certificação de engate coaxial dependeria de mecanismos adicionais, como guias mecânicas, cones de acoplamento ou estratégias de controle específicas para a fase de contato.

\subsection{Encerramento e limitações do modelo}

Os resultados indicam que a arquitetura proposta para a missão de \gls{rvdb} atinge os objetivos estabelecidos: a transferência orbital gera condições adequadas para o uso do referencial LVLH; o controle translacional baseado nas equações de \gls{hcw} conduz o perseguidor à vizinhança do alvo com limites de posição e velocidade; o controle de atitude mantém o alinhamento com o alvo mesmo com o manipulador em operação; e o manipulador é capaz de posicionar a ferramenta em uma vizinhança milimétrica da posição alvo, com baixo esforço de controle e efeitos limitados sobre a base.

Por outro lado, o modelo apresenta limitações que devem ser consideradas na interpretação dos resultados. A dinâmica relativa é descrita por equações linearizadas \gls{hcw} em órbitas circulares coplanares, sem consideração explícita de perturbações orbitais, atrasos de atuadores, ruídos de sensores ou saturações detalhadas de torque. O manipulador é representado por um modelo rígido de três graus de liberdade, sem flexibilidade estrutural nem dinâmica de contato com a porta de acoplamento. Além disso, a condição de alinhamento angular entre ferramenta e porta não possui modelagem da interface mecânica de engate.

% Embora hajam simplificações em termos de alinhamento com a porta de acoplamento, a coerência entre as fases da missão, os níveis de esforço de controle obtidos, os baixos valores residuais de energia e momento angular e a precisão geométrica alcançada na posição da ferramenta indicam que a estratégia de controle adotada é compatível com o cenário estudado e fornece base consistente para estudos futuros de extensões do modelo, incluindo dinâmica de contato, modelos orbitais mais completos e técnicas de controle alternativas para a fase de aproximação final.

